\documentclass[11pt]{article}

\usepackage[margin=1in]{geometry}
\usepackage{amsmath,amssymb,bm}
\usepackage{booktabs}
\usepackage{array}
\usepackage[hidelinks]{hyperref}

% ---------------------------------------------------------------------------
% Notation macros (shared with solid_mechanics_forms.tex where they overlap)
% ---------------------------------------------------------------------------
\newcommand{\ud}{\,\mathrm{d}}
\newcommand{\Grad}{\operatorname{Grad}}
\newcommand{\Dive}{\operatorname{Div}}
\newcommand{\dev}{\operatorname{dev}}
\newcommand{\tr}{\operatorname{tr}}
\newcommand{\cof}{\operatorname{cof}}
\newcommand{\bx}{\bm{x}}
\newcommand{\bX}{\bm{X}}
\newcommand{\bu}{\bm{u}}
\newcommand{\bw}{\bm{w}}
\newcommand{\be}{\bm{e}}
\newcommand{\bn}{\bm{n}}
\newcommand{\bN}{\bm{N}}
\newcommand{\bt}{\bm{t}}
\newcommand{\bsig}{\bm{\sigma}}
\newcommand{\beps}{\bm{\varepsilon}}
\newcommand{\bP}{\bm{P}}
\newcommand{\bS}{\bm{S}}
\newcommand{\bF}{\bm{F}}
\newcommand{\bFb}{\bar{\bm{F}}}
\newcommand{\bC}{\bm{C}}
\newcommand{\bB}{\bm{B}}
\newcommand{\bI}{\bm{I}}
\newcommand{\bphi}{\bm{\varphi}}
\newcommand{\Ione}{I_1}
\newcommand{\Itwo}{I_2}
\newcommand{\Ib}{\bar{I}_1}
\newcommand{\IIb}{\bar{I}_2}
\newcommand{\Pone}{\Psi_{1}}
\newcommand{\Ptwo}{\Psi_{2}}
\newcommand{\code}[1]{{\frenchspacing\texttt{#1}}}

\title{Incompressible Isotropic Hyperelasticity:\\
Constitutive Models and Homogeneous Solutions}
\author{}
\date{}

\begin{document}
\maketitle

\noindent
This note collects the standard incompressible isotropic hyperelastic
models and the closed-form stress states they predict under homogeneous
deformation (uniaxial, equibiaxial, pure shear, simple shear).  The
homogeneous solutions serve as reference solutions for the mixed
displacement--pressure formulation of the framework
(\code{formulation: mixed} with any of the six models at $\kappa = \infty$;
Section \ref{sec:verification}).  Conventions follow the code:
the pressure unknown $p$ is the tensile-positive mean stress, and in two
dimensions the kernels are plane strain with $F_{33} = 1$ unless plane
stress is selected (Section \ref{sec:planestress}).  The weak forms
themselves are in \code{solid\_mechanics\_forms.tex}.

%============================================================================
\section{Kinematics and the Incompressibility Constraint}
%============================================================================

Let $\Omega_R$ be the reference configuration; referential quantities
carry the subscript $R$.  With $\bx = \bphi(\bX) = \bX + \bu(\bX)$,
$\bX \in \Omega_R$, the deformation gradient, Jacobian, right and left
Cauchy--Green tensors are
\begin{equation}
  \bF = \bI + \Grad\bu, \qquad J = \det\bF, \qquad
  \bC = \bF^T\bF, \qquad \bB = \bF\bF^T .
\end{equation}
The principal invariants of $\bC$ (equivalently of $\bB$) are
\begin{equation}
  \Ione = \tr\bC, \qquad
  \Itwo = \tfrac12\bigl[(\tr\bC)^2 - \tr(\bC^2)\bigr] = \tr(\cof\bC), \qquad
  I_3 = \det\bC = J^2 .
\end{equation}
In terms of the principal stretches $\lambda_i$ (the square roots of the
eigenvalues of $\bC$ and $\bB$),
\begin{equation}
  \Ione = \lambda_1^2 + \lambda_2^2 + \lambda_3^2, \qquad
  \Itwo = \lambda_1^2\lambda_2^2 + \lambda_2^2\lambda_3^2 + \lambda_3^2\lambda_1^2, \qquad
  J = \lambda_1\lambda_2\lambda_3 .
\end{equation}
An incompressible material satisfies the internal constraint
\begin{equation}
  J = 1 \qquad (\lambda_1\lambda_2\lambda_3 = 1),
  \label{eq:constraint}
\end{equation}
under which two identities used repeatedly below hold:
\begin{equation}
  \Itwo = \lambda_1^{-2} + \lambda_2^{-2} + \lambda_3^{-2} = \tr\bB^{-1},
  \qquad
  \bB^2 = \Ione\,\bB - \Itwo\,\bI + \bB^{-1}
  \quad\text{(Cayley--Hamilton with } I_3 = 1\text{)}.
  \label{eq:ch}
\end{equation}

\paragraph{Isochoric split.}
The framework writes its decoupled models in terms of the isochoric
deformation gradient and modified invariants
\begin{equation}
  \bFb = J^{-1/3}\bF, \qquad \Ib = J^{-2/3}\Ione, \qquad \IIb = J^{-4/3}\Itwo,
\end{equation}
so that the same energy makes sense at finite bulk modulus.  On an
incompressible solution $J = 1$ and the barred and unbarred quantities
coincide; the formulas below are written with $\Ione, \Itwo$.

%============================================================================
\section{Constitutive Framework}
%============================================================================

\subsection{Strain energy and the reaction stress}

For an isotropic incompressible material the strain energy per unit
reference volume is a function of the two remaining invariants,
$\Psi = \Psi(\Ione,\Itwo)$, or equivalently a symmetric function of the
principal stretches subject to \eqref{eq:constraint}.  Write
\begin{equation}
  \Pone = \frac{\partial\Psi}{\partial\Ione}, \qquad
  \Ptwo = \frac{\partial\Psi}{\partial\Itwo}.
\end{equation}
The constraint is enforced by a Lagrange multiplier $q$ that does no work,
and the Cauchy stress is
\begin{equation}
  \bsig = 2\Pone\,\bB + 2\Ptwo\,(\Ione\bB - \bB^2) - q\,\bI
        = 2\Pone\,\bB - 2\Ptwo\,\bB^{-1} + (2\Ptwo\Itwo - q)\,\bI ,
  \label{eq:sigma_q}
\end{equation}
where the second form uses \eqref{eq:ch}.  The multiplier $q$ is not given
by the constitutive law; it is determined by equilibrium and the boundary
conditions.  Only the isotropic part of $\bsig$ is affected, so the
convention-free statement of \eqref{eq:sigma_q} is the deviatoric--mean
split
\begin{equation}
  \boxed{\;
  \bsig = \bsig' + p\,\bI, \qquad
  \bsig' = 2\Pone\,\dev\bB - 2\Ptwo\,\dev\bB^{-1}, \qquad
  p = \tfrac13\tr\bsig ,\;}
  \label{eq:sigma_split}
\end{equation}
with $\tr\bB = \Ione$ and $\tr\bB^{-1} = \Itwo$ in the deviators.

\paragraph{Sign convention.}
Much of the literature writes $\bsig = -\tilde p\,\bI + 2\Pone\bB - 2\Ptwo\bB^{-1}$
and calls $\tilde p$ the hydrostatic pressure.  Because different authors
absorb different isotropic terms into $\tilde p$, its value is
convention-dependent; $p = \tfrac13\tr\bsig$ (tensile positive) is not, and
it is the quantity the mixed formulation solves for (Section
\ref{sec:mixed}).  All results below are stated in terms of $p$.

\subsection{Principal-stretch form}

In the principal frame $\bB = \sum_i\lambda_i^2\,\bn_i\otimes\bn_i$ and
$\bsig = \sum_i\sigma_i\,\bn_i\otimes\bn_i$ with
\begin{equation}
  \sigma_i = 2\Pone\lambda_i^2 - 2\Ptwo\lambda_i^{-2} - q
           = \lambda_i\frac{\partial\Psi}{\partial\lambda_i} - q' ,
  \label{eq:principal}
\end{equation}
where $q' = q - 2\Ptwo\Itwo$ and the second form follows from
$\partial\Ione/\partial\lambda_i = 2\lambda_i$,
$\partial\Itwo/\partial\lambda_i = 2\lambda_i(\Ione - \lambda_i^2)$ together
with \eqref{eq:constraint}:
\begin{equation}
  \lambda_i\frac{\partial\Psi}{\partial\lambda_i}
  = 2\Pone\lambda_i^2 - 2\Ptwo\lambda_i^{-2} + 2\Ptwo\Itwo .
\end{equation}
The principal stress differences are constitutive,
\begin{equation}
  \sigma_i - \sigma_j
  = 2\Pone(\lambda_i^2 - \lambda_j^2) - 2\Ptwo(\lambda_i^{-2} - \lambda_j^{-2})
  = \lambda_i\frac{\partial\Psi}{\partial\lambda_i}
  - \lambda_j\frac{\partial\Psi}{\partial\lambda_j},
  \label{eq:differences}
\end{equation}
and the recipe for every homogeneous problem is the same: the two
differences follow from \eqref{eq:differences}, one principal stress is
fixed by a traction-free face, and then $p = \tfrac13(\sigma_1 + \sigma_2 + \sigma_3)$.

\subsection{Other stress measures}

With $J = 1$ the first and second Piola--Kirchhoff stresses are
\begin{equation}
  \bP = \bsig\bF^{-T}, \qquad \bS = \bF^{-1}\bsig\bF^{-T},
\end{equation}
and for a diagonal $\bF = \operatorname{diag}(\lambda_1,\lambda_2,\lambda_3)$
\begin{equation}
  P_i = \frac{\sigma_i}{\lambda_i}, \qquad S_i = \frac{\sigma_i}{\lambda_i^2}.
\end{equation}
$P_i$ is the nominal stress (force per unit reference area), which is what
a traction condition of the total Lagrangian formulation prescribes:
$\bt_R = \bP\bN$ on the reference boundary $\partial\Omega_R$.

\subsection{Small-strain limit}

For $\bu = \beps$-small, $\bB \approx \bI + 2\beps$,
$\bB^{-1} \approx \bI - 2\beps$ and $\tr\beps = 0$, so
\eqref{eq:sigma_split} gives
\begin{equation}
  \bsig \approx 2\mu\,\beps + p\,\bI, \qquad
  \mu = 2\bigl(\Pone + \Ptwo\bigr)\big|_{\Ione = \Itwo = 3},
  \label{eq:mu}
\end{equation}
which identifies the small-strain shear modulus of every model.  The
corresponding Young's modulus and Poisson ratio are $E = 3\mu$, $\nu = 1/2$.

\subsection{Relation to the mixed formulation}
\label{sec:mixed}

The framework's decoupled materials provide $\Psi_{\mathrm{iso}}(\bFb)$
and $\bP_{\mathrm{iso}} = \partial\Psi_{\mathrm{iso}}/\partial\bF$.  Since
$\Psi_{\mathrm{iso}}$ is invariant under $\bF \to \alpha\bF$, Euler's
theorem gives $\bP_{\mathrm{iso}}:\bF = 0$, i.e.\
$\tr(J^{-1}\bP_{\mathrm{iso}}\bF^T) = 0$: the isochoric Cauchy stress is
deviatoric for every $\bF$, not only at $J = 1$.  The mixed kernel uses
\begin{equation}
  \bP = \bP_{\mathrm{iso}}(\bF) + p\,J\bF^{-T}
  \quad\Longrightarrow\quad
  \bsig = J^{-1}\bP\bF^T = \underbrace{J^{-1}\bP_{\mathrm{iso}}\bF^T}_{\text{deviatoric}} + p\,\bI,
\end{equation}
so its pressure unknown is exactly $p = \tfrac13\tr\bsig$ of
\eqref{eq:sigma_split}, and on an incompressible solution
$J^{-1}\bP_{\mathrm{iso}}\bF^T = \bsig'$.  The constraint equation is
$\int q\,(J - 1 - p/\kappa)\ud V = 0$: for $\kappa = \infty$ it is
\eqref{eq:constraint} in weak form and $p$ is a Lagrange multiplier; for
finite $\kappa$ it is the penalty relation $p = \kappa(J - 1)$ and the
solutions below are recovered as $\kappa \to \infty$.  (Homogeneous
solutions at finite $\kappa$ require solving a scalar nonlinear equation
for the lateral stretch and are not tabulated here.)

%============================================================================
\section{Constitutive Models}
%============================================================================

Each model is given by $\Psi$, its derivatives $\Pone,\Ptwo$ (or
$\lambda_i\partial\Psi/\partial\lambda_i$), and its small-strain shear
modulus from \eqref{eq:mu}.  All are normalised so that $\Psi = 0$ and
$\bsig = p\bI$ in the reference configuration.

\subsection{Neo-Hookean}
\begin{equation}
  \Psi = \frac{\mu}{2}(\Ione - 3), \qquad
  \Pone = \frac{\mu}{2}, \quad \Ptwo = 0, \qquad
  \bsig = \mu\,\dev\bB + p\,\bI .
\end{equation}
One parameter, the shear modulus $\mu$.  This is the Gaussian-chain model of
rubber elasticity (Treloar) and the incompressible limit of
\code{iso\_neo\_hookean}.

\subsection{Mooney--Rivlin}
\begin{gather}
  \Psi = c_1(\Ione - 3) + c_2(\Itwo - 3), \qquad
  \Pone = c_1, \quad \Ptwo = c_2, \qquad \mu = 2(c_1 + c_2), \\
  \bsig = 2c_1\,\dev\bB - 2c_2\,\dev\bB^{-1} + p\,\bI .
\end{gather}
The $c_2$ term is what distinguishes the model from neo-Hookean: it
changes the uniaxial and equibiaxial curves and the stresses in the
constrained directions, whereas the pure-shear stress $\sigma_1$ and the
simple-shear stress $\sigma_{12}$ depend on $\mu$ alone (Section
\ref{sec:homogeneous}).  This is \code{mooney\_rivlin} in the code.  The general Rivlin polynomial
\begin{equation}
  \Psi = \sum_{i+j \ge 1} C_{ij}\,(\Ione - 3)^i(\Itwo - 3)^j,
  \qquad \mu = 2(C_{10} + C_{01}),
\end{equation}
contains neo-Hookean ($C_{10} = \mu/2$), Mooney--Rivlin
($C_{10} = c_1$, $C_{01} = c_2$) and Yeoh as special cases.

\subsection{Yeoh}
\begin{equation}
  \Psi = C_{10}(\Ione - 3) + C_{20}(\Ione - 3)^2 + C_{30}(\Ione - 3)^3, \qquad
  \Pone = C_{10} + 2C_{20}(\Ione - 3) + 3C_{30}(\Ione - 3)^2, \quad \Ptwo = 0,
\end{equation}
with $\mu = 2C_{10}$.  A cubic in $\Ione$ alone; typically $C_{20} < 0$ and
$C_{30} > 0$, which reproduces the softening-then-stiffening uniaxial
response of filled rubbers.  This is \code{yeoh} in the code.

\subsection{Gent}
\begin{equation}
  \Psi = -\frac{\mu J_m}{2}\ln\!\Bigl(1 - \frac{\Ione - 3}{J_m}\Bigr), \qquad
  \Pone = \frac{\mu}{2}\,\frac{J_m}{J_m - (\Ione - 3)}, \quad \Ptwo = 0,
  \qquad \Ione - 3 < J_m .
\end{equation}
Two parameters, $\mu$ and the limiting value $J_m$ of $\Ione - 3$
(limited chain extensibility).  $J_m \to \infty$ recovers neo-Hookean.  This
is \code{gent} in the code.

\subsection{Arruda--Boyce (eight-chain)}
In the usual five-term series form,
\begin{equation}
  \Psi = \mu\sum_{i=1}^{5}\frac{C_i}{N^{\,i-1}}\bigl(\Ione^i - 3^i\bigr), \qquad
  C_1 = \tfrac12,\ C_2 = \tfrac{1}{20},\ C_3 = \tfrac{11}{1050},\
  C_4 = \tfrac{19}{7000},\ C_5 = \tfrac{519}{673750},
\end{equation}
\begin{equation}
  \Pone = \mu\sum_{i=1}^{5}\frac{i\,C_i}{N^{\,i-1}}\,\Ione^{\,i-1}
        = \mu\Bigl[\tfrac12 + \tfrac{\Ione}{10N} + \tfrac{11\Ione^2}{350N^2}
          + \tfrac{19\Ione^3}{1750N^3} + \tfrac{519\Ione^4}{134750N^4}\Bigr],
  \qquad \Ptwo = 0 .
\end{equation}
Here $\mu$ is the chain-density parameter $nkT$ and $N$ the number of
rigid links per chain (the locking stretch is $\lambda_L = \sqrt N$).  The
small-strain shear modulus is not $\mu$ but
\begin{equation}
  \mu_0 = 2\Pone\big|_{\Ione = 3}
  = \mu\Bigl(1 + \frac{3}{5N} + \frac{99}{175N^2} + \frac{513}{875N^3} + \frac{42039}{67375N^4}\Bigr).
  \label{eq:ab_mu0}
\end{equation}
The series form is \code{arruda\_boyce} in the code.  The exact form uses the
inverse Langevin function
$\mathcal L(x) = \coth x - 1/x$: with $\lambda_{ch} = \sqrt{\Ione/3}$ and
$\beta = \mathcal L^{-1}(\lambda_{ch}/\sqrt N)$,
\begin{equation}
  \Psi = \mu N\Bigl[\frac{\lambda_{ch}}{\sqrt N}\,\beta + \ln\frac{\beta}{\sinh\beta}\Bigr] + \text{const},
  \qquad
  \Pone = \frac{\mu}{6}\,\frac{\sqrt N}{\lambda_{ch}}\,\mathcal L^{-1}\!\Bigl(\frac{\lambda_{ch}}{\sqrt N}\Bigr).
\end{equation}

\subsection{Ogden}
Expressed directly in the principal stretches,
\begin{equation}
  \Psi = \sum_{r=1}^{M}\frac{\mu_r}{\alpha_r}
         \bigl(\lambda_1^{\alpha_r} + \lambda_2^{\alpha_r} + \lambda_3^{\alpha_r} - 3\bigr),
  \qquad
  \lambda_i\frac{\partial\Psi}{\partial\lambda_i} = \sum_{r}\mu_r\lambda_i^{\alpha_r},
  \qquad
  \mu = \frac12\sum_r\mu_r\alpha_r ,
\end{equation}
with $\mu_r\alpha_r > 0$ for each term.  The exponents need not be
integers.  Neo-Hookean is $M = 1$, $\alpha_1 = 2$, $\mu_1 = \mu$;
Mooney--Rivlin is $M = 2$, $(\alpha_1,\alpha_2) = (2,-2)$,
$(\mu_1,\mu_2) = (2c_1, -2c_2)$, using $\sum\lambda_i^{-2} = \Itwo$.  Ogden's
classical three-term fit to Treloar's vulcanised-rubber data is
$(\mu_r) = (0.63,\ 0.0012,\ -0.01)\,\mathrm{MPa}$,
$(\alpha_r) = (1.3,\ 5.0,\ -2.0)$.  Because it is not an invariant model,
Ogden's stresses are evaluated in the principal frame with
\eqref{eq:principal}--\eqref{eq:differences}.  This is \code{ogden} in the
code, which evaluates it spectrally (Jacobi eigen-decomposition of $\bC$)
and forms the tangent analytically from the spectral representation, with
the usual limit at coincident principal stretches, since eigenvectors are
not differentiable there.

\begin{table}[htbp]
\centering
\small
\begin{tabular}{@{}l l l l@{}}
\toprule
Model & $\Psi$ & Parameters & $\mu$ \\
\midrule
Neo-Hookean & $\frac{\mu}{2}(\Ione-3)$ & $\mu$ & $\mu$ \\
Mooney--Rivlin & $c_1(\Ione-3)+c_2(\Itwo-3)$ & $c_1, c_2$ & $2(c_1+c_2)$ \\
Yeoh & $\sum_{i=1}^{3}C_{i0}(\Ione-3)^i$ & $C_{10},C_{20},C_{30}$ & $2C_{10}$ \\
Gent & $-\frac{\mu J_m}{2}\ln\bigl(1-\frac{\Ione-3}{J_m}\bigr)$ & $\mu, J_m$ & $\mu$ \\
Arruda--Boyce & $\mu\sum_{i=1}^{5}C_iN^{1-i}(\Ione^i-3^i)$ & $\mu, N$ & $\mu(1+\frac{3}{5N}+\dots)$ \\
Ogden & $\sum_r\frac{\mu_r}{\alpha_r}(\sum_i\lambda_i^{\alpha_r}-3)$ & $\mu_r,\alpha_r$ & $\frac12\sum_r\mu_r\alpha_r$ \\
\bottomrule
\end{tabular}
\caption{Incompressible isotropic models and their small-strain shear modulus.}
\label{tab:models}
\end{table}

\paragraph{$\Ione$-only models.}
Yeoh, Gent and Arruda--Boyce depend on $\Ione$ alone, so $\Ptwo = 0$ and
\eqref{eq:sigma_split} reduces to the neo-Hookean form with a
deformation-dependent modulus,
\begin{equation}
  \bsig = \mu_{\mathrm{eff}}(\Ione)\,\dev\bB + p\,\bI, \qquad
  \mu_{\mathrm{eff}}(\Ione) = 2\Pone(\Ione).
  \label{eq:mueff}
\end{equation}
Every neo-Hookean result in Section \ref{sec:homogeneous} therefore
applies to these models after replacing $\mu$ by
$\mu_{\mathrm{eff}}$ evaluated at the $\Ione$ of that deformation.

%============================================================================
\section{Homogeneous Deformations}
\label{sec:homogeneous}
%============================================================================

A homogeneous deformation has constant $\bF$, hence constant $\bsig$ and
$\bP$; $\Dive\bP = 0$ holds identically, so the state is an exact solution
of the equilibrium equations without body force, and $p$ is a constant
fixed by the traction-free faces.  For each mode below the results are
given for a general $\Psi(\Ione,\Itwo)$, then specialised to neo-Hookean
(NH), Mooney--Rivlin (MR) and Ogden; $\Ione$-only models use
\eqref{eq:mueff}.  Extension is $\lambda > 1$, compression $\lambda < 1$;
the formulas cover both.

\subsection{Uniaxial tension and compression}
\begin{equation}
  \bF = \operatorname{diag}(\lambda,\ \lambda^{-1/2},\ \lambda^{-1/2}), \qquad
  \Ione = \lambda^2 + \frac{2}{\lambda}, \qquad
  \Itwo = 2\lambda + \frac{1}{\lambda^2}, \qquad
  \sigma_2 = \sigma_3 = 0 .
\end{equation}
From \eqref{eq:differences} with $\lambda_1 = \lambda$,
$\lambda_2 = \lambda^{-1/2}$,
\begin{equation}
  \boxed{\;
  \sigma_1 = 2\Bigl(\lambda^2 - \frac1\lambda\Bigr)\Bigl(\Pone + \frac{\Ptwo}{\lambda}\Bigr),
  \qquad
  P_1 = \frac{\sigma_1}{\lambda} = 2\Bigl(\lambda - \frac1{\lambda^2}\Bigr)\Bigl(\Pone + \frac{\Ptwo}{\lambda}\Bigr),
  \qquad
  p = \frac{\sigma_1}{3}.\;}
\end{equation}
\begin{align}
  \text{NH:}\quad & \sigma_1 = \mu\Bigl(\lambda^2 - \frac1\lambda\Bigr), &
  P_1 &= \mu\Bigl(\lambda - \frac1{\lambda^2}\Bigr), \\
  \text{MR:}\quad & \sigma_1 = 2\Bigl(c_1 + \frac{c_2}{\lambda}\Bigr)\Bigl(\lambda^2 - \frac1\lambda\Bigr), &
  P_1 &= 2\Bigl(c_1 + \frac{c_2}{\lambda}\Bigr)\Bigl(\lambda - \frac1{\lambda^2}\Bigr), \\
  \text{Ogden:}\quad & \sigma_1 = \sum_r\mu_r\bigl(\lambda^{\alpha_r} - \lambda^{-\alpha_r/2}\bigr), &
  P_1 &= \sum_r\mu_r\bigl(\lambda^{\alpha_r-1} - \lambda^{-\alpha_r/2-1}\bigr).
\end{align}
Small strain: $\sigma_1 \approx 3\mu(\lambda - 1) = E(\lambda - 1)$.

\subsection{Equibiaxial tension}
\begin{equation}
  \bF = \operatorname{diag}(\lambda,\ \lambda,\ \lambda^{-2}), \qquad
  \Ione = 2\lambda^2 + \frac1{\lambda^4}, \qquad
  \Itwo = \lambda^4 + \frac{2}{\lambda^2}, \qquad
  \sigma_3 = 0 .
\end{equation}
\begin{equation}
  \boxed{\;
  \sigma_1 = \sigma_2 = 2\Bigl(\lambda^2 - \frac1{\lambda^4}\Bigr)\bigl(\Pone + \lambda^2\Ptwo\bigr),
  \qquad
  P_1 = P_2 = \frac{\sigma_1}{\lambda},
  \qquad
  p = \frac{2\sigma_1}{3}.\;}
\end{equation}
\begin{align}
  \text{NH:}\quad & \sigma_1 = \mu\Bigl(\lambda^2 - \frac1{\lambda^4}\Bigr), \\
  \text{MR:}\quad & \sigma_1 = 2\bigl(c_1 + c_2\lambda^2\bigr)\Bigl(\lambda^2 - \frac1{\lambda^4}\Bigr), \\
  \text{Ogden:}\quad & \sigma_1 = \sum_r\mu_r\bigl(\lambda^{\alpha_r} - \lambda^{-2\alpha_r}\bigr).
\end{align}
Equibiaxial tension with stretch $\lambda$ is uniaxial compression to
$\lambda^{-2}$ plus a hydrostatic stress, which is why the two tests are
interchangeable for identifying $\Psi$.

\subsection{Pure shear (planar tension)}
\begin{equation}
  \bF = \operatorname{diag}(\lambda,\ 1,\ \lambda^{-1}), \qquad
  \Ione = \Itwo = \lambda^2 + 1 + \frac1{\lambda^2}, \qquad
  \sigma_3 = 0 .
\end{equation}
\begin{equation}
  \boxed{\;
  \begin{aligned}
  \sigma_1 &= 2\Bigl(\lambda^2 - \frac1{\lambda^2}\Bigr)\bigl(\Pone + \Ptwo\bigr), &
  \sigma_2 &= 2\Bigl(1 - \frac1{\lambda^2}\Bigr)\bigl(\Pone + \lambda^2\Ptwo\bigr), \\
  P_1 &= \frac{\sigma_1}{\lambda},\quad P_2 = \sigma_2, &
  p &= \frac{\sigma_1 + \sigma_2}{3}.
  \end{aligned}\;}
  \label{eq:pureshear}
\end{equation}
\begin{align}
  \text{NH:}\quad & \sigma_1 = \mu\Bigl(\lambda^2 - \frac1{\lambda^2}\Bigr), &
  \sigma_2 &= \mu\Bigl(1 - \frac1{\lambda^2}\Bigr), \\
  \text{MR:}\quad & \sigma_1 = \mu\Bigl(\lambda^2 - \frac1{\lambda^2}\Bigr),\ \ \mu = 2(c_1 + c_2), &
  \sigma_2 &= 2\bigl(c_1 + c_2\lambda^2\bigr)\Bigl(1 - \frac1{\lambda^2}\Bigr), \\
  \text{Ogden:}\quad & \sigma_1 = \sum_r\mu_r\bigl(\lambda^{\alpha_r} - \lambda^{-\alpha_r}\bigr), &
  \sigma_2 &= \sum_r\mu_r\bigl(1 - \lambda^{-\alpha_r}\bigr).
\end{align}
The name refers to the principal stretches $(\lambda, 1, \lambda^{-1})$,
which are those of a simple shear of amount $\gamma = \lambda - \lambda^{-1}$;
pure shear is simple shear followed by a rotation.  Note that for
Mooney--Rivlin $\sigma_1$ depends only on $\mu$; the constrained-direction
stress $\sigma_2$ is what distinguishes $c_1$ from $c_2$.

\subsection{Simple shear}
\begin{equation}
  \bF = \bI + \gamma\,\be_1\otimes\be_2
  \quad (x_1 = X_1 + \gamma X_2), \quad
  \bB = \begin{pmatrix} 1+\gamma^2 & \gamma & 0 \\ \gamma & 1 & 0 \\ 0 & 0 & 1\end{pmatrix}, \quad
  \bB^{-1} = \begin{pmatrix} 1 & -\gamma & 0 \\ -\gamma & 1+\gamma^2 & 0 \\ 0 & 0 & 1\end{pmatrix},
\end{equation}
with $\Ione = \Itwo = 3 + \gamma^2$.  The principal stretches are
$\lambda = \tfrac12\bigl(\gamma + \sqrt{\gamma^2 + 4}\bigr)$, $\lambda^{-1}$ and
$1$ (so $\gamma = \lambda - \lambda^{-1}$), with the in-plane principal
directions at an angle $\theta$ to $\be_1$ given by $\tan 2\theta = 2/\gamma$.
The deviatoric stress \eqref{eq:sigma_split} is
\begin{equation}
  \sigma'_{12} = 2(\Pone + \Ptwo)\,\gamma, \qquad
  \sigma'_{11} = \tfrac23(2\Pone + \Ptwo)\gamma^2, \qquad
  \sigma'_{22} = -\tfrac23(\Pone + 2\Ptwo)\gamma^2, \qquad
  \sigma'_{33} = -\tfrac23(\Pone - \Ptwo)\gamma^2,
\end{equation}
which gives the shear stress and the normal-stress differences
independently of $p$:
\begin{equation}
  \boxed{\;
  \begin{aligned}
  \sigma_{12} &= 2(\Pone + \Ptwo)\,\gamma, &
  \sigma_{11} - \sigma_{22} &= \gamma\,\sigma_{12} = 2(\Pone + \Ptwo)\gamma^2, \\
  \sigma_{11} - \sigma_{33} &= 2\Pone\gamma^2, &
  \sigma_{22} - \sigma_{33} &= -2\Ptwo\gamma^2 .
  \end{aligned}\;}
\end{equation}
The first relation defines the secant shear modulus
$\mu(\gamma) = \sigma_{12}/\gamma = 2(\Pone + \Ptwo)$, which is constant for
neo-Hookean and Mooney--Rivlin (both predict a linear shear response) and
depends on $\gamma$ through $\Ione = 3 + \gamma^2$ for the $\Ione$-only
models.  The relation $\sigma_{11} - \sigma_{22} = \gamma\sigma_{12}$ is
universal (independent of $\Psi$).  The mean stress depends on which face
is traction-free:
\begin{align}
  \sigma_{33} = 0\ \text{(thin sheet):}\quad &
  \sigma_{11} = 2\Pone\gamma^2, & \sigma_{22} &= -2\Ptwo\gamma^2, &
  p &= \tfrac23(\Pone - \Ptwo)\gamma^2, \\
  \sigma_{22} = 0\ \text{(free top face):}\quad &
  \sigma_{11} = 2(\Pone + \Ptwo)\gamma^2, & \sigma_{33} &= 2\Ptwo\gamma^2, &
  p &= \tfrac23(\Pone + 2\Ptwo)\gamma^2 .
\end{align}
The second case is the plane-strain one.
Specialised:
\begin{align}
  \text{NH (either case):}\quad & \sigma_{12} = \mu\gamma, \quad \sigma_{11} = \mu\gamma^2, \quad
  \sigma_{22} = \sigma_{33} = 0, \quad p = \frac{\mu\gamma^2}{3}, \\
  \text{MR, } \sigma_{33} = 0:\quad & \sigma_{12} = \mu\gamma, \quad \sigma_{11} = 2c_1\gamma^2, \quad
  \sigma_{22} = -2c_2\gamma^2, \quad p = \frac{2(c_1 - c_2)\gamma^2}{3}, \\
  \text{MR, } \sigma_{22} = 0:\quad & \sigma_{12} = \mu\gamma, \quad \sigma_{11} = \mu\gamma^2, \quad
  \sigma_{33} = 2c_2\gamma^2, \quad p = \frac{2(c_1 + 2c_2)\gamma^2}{3}, \\
  \text{Ogden:}\quad & \sigma_{12} = \frac{\sum_r\mu_r(\lambda^{\alpha_r} - \lambda^{-\alpha_r})}{\lambda + \lambda^{-1}}, \quad
  \sigma_{11} - \sigma_{22} = \gamma\sigma_{12}, \notag \\
  & \sigma_{11} + \sigma_{22} - 2\sigma_{33} = \sum_r\mu_r\bigl(\lambda^{\alpha_r} + \lambda^{-\alpha_r} - 2\bigr).
\end{align}
The nominal stress follows from $\bP = \bsig\bF^{-T}$ with
$\bF^{-T} = \bI - \gamma\,\be_2\otimes\be_1$:
\begin{equation}
  P_{11} = \sigma_{11} - \gamma\sigma_{12}, \qquad
  P_{12} = \sigma_{12}, \qquad
  P_{21} = \sigma_{12} - \gamma\sigma_{22}, \qquad
  P_{22} = \sigma_{22}, \qquad
  P_{33} = \sigma_{33},
\end{equation}
so that for neo-Hookean $P_{11} = P_{22} = 0$ and $P_{12} = P_{21} = \mu\gamma$:
the nominal tractions on the sheared faces are purely tangential.  For any
other model the lateral faces $X_1 = \text{const}$ must carry the normal
nominal traction $P_{11}$ for the simple shear to be an exact solution.

\subsection{Consistency checks}
All stresses vanish at $\lambda = 1$ or $\gamma = 0$.  The initial slopes,
with $\mu$ from \eqref{eq:mu}, are
\begin{equation}
  \frac{\ud\sigma_1}{\ud\lambda}\Big|_{\lambda=1} =
  \begin{cases} 3\mu & \text{uniaxial} \\ 6\mu & \text{equibiaxial} \\ 4\mu & \text{pure shear}, \end{cases}
  \qquad
  \frac{\ud\sigma_{12}}{\ud\gamma}\Big|_{\gamma=0} = \mu ,
\end{equation}
and the normal-stress relation $\sigma_{11} - \sigma_{22} = \gamma\sigma_{12}$
in simple shear holds for every isotropic incompressible model.

\begin{table}[htbp]
\centering
\small
\renewcommand{\arraystretch}{1.5}
\begin{tabular}{@{}l l l l l@{}}
\toprule
Mode & Stretches & $\Ione$ & Non-zero Cauchy stress & $p = \tfrac13\tr\bsig$ \\
\midrule
Uniaxial & $(\lambda, \lambda^{-1/2}, \lambda^{-1/2})$ & $\lambda^2 + 2\lambda^{-1}$ &
  $\sigma_1 = \mu(\lambda^2 - \lambda^{-1})$ & $\sigma_1/3$ \\
Equibiaxial & $(\lambda, \lambda, \lambda^{-2})$ & $2\lambda^2 + \lambda^{-4}$ &
  $\sigma_1 = \sigma_2 = \mu(\lambda^2 - \lambda^{-4})$ & $2\sigma_1/3$ \\
Pure shear & $(\lambda, 1, \lambda^{-1})$ & $\lambda^2 + 1 + \lambda^{-2}$ &
  $\sigma_1 = \mu(\lambda^2 - \lambda^{-2}),\ \sigma_2 = \mu(1 - \lambda^{-2})$ & $(\sigma_1 + \sigma_2)/3$ \\
Simple shear & $(\lambda, \lambda^{-1}, 1)$, $\gamma = \lambda - \lambda^{-1}$ & $3 + \gamma^2$ &
  $\sigma_{12} = \mu\gamma,\ \sigma_{11} = \mu\gamma^2$ & $\mu\gamma^2/3$ \\
\bottomrule
\end{tabular}
\caption{Neo-Hookean homogeneous solutions (one face traction-free as
stated in the text).  For Yeoh, Gent and Arruda--Boyce replace $\mu$ by
$2\Pone(\Ione)$ at the tabulated $\Ione$.}
\label{tab:nh}
\end{table}

%============================================================================
\section{Plane Strain and Reproducing These States in the Code}
%============================================================================

In two dimensions the kernels build $\bF = \begin{pmatrix}\bF_{2D} & 0\\ 0 & 1\end{pmatrix}$
and $J = \det\bF_{2D}$, i.e.\ plane strain with $\lambda_3 = 1$ (the default,
\code{plane: strain}; Section \ref{sec:planestress} for plane stress).  An
incompressible plane-strain state therefore has in-plane principal
stretches $(\lambda, \lambda^{-1})$: up to a rotation every such
homogeneous state is a pure shear \eqref{eq:pureshear} with $\be_3$ as the
constrained direction.  The out-of-plane stress $\sigma_{33}$ is a reaction
and is generally non-zero; it enters both the pressure field
$p = \tfrac13(\sigma_{11} + \sigma_{22} + \sigma_{33})$ and the
\code{vonmises} output, which is computed from the full $3\times3$ Cauchy
stress, $\sigma_{vm} = \sqrt{\tfrac32\,\bsig':\bsig'}$.  The code can also
write $\bsig$ (six components), $\bP$ and $\bF$ (nine) and the energy
density themselves.  All such quantities are computed at the quadrature
points and then presented as a continuous nodal field, as element averages,
or as the raw point cloud (\code{output.fields}, \code{output.quadrature\_at},
\code{output.nodal\_projection}).

\subsection{Plane-strain extension}
Stretch $\lambda$ along $\be_1$ with the faces $X_2 = 0, H$ traction-free
($\sigma_{22} = 0$).  Relabelling \eqref{eq:pureshear} with
$(\lambda_1,\lambda_2,\lambda_3) = (\lambda, \lambda^{-1}, 1)$:
\begin{equation}
  \bu = (\lambda - 1)X_1\,\be_1 + (\lambda^{-1} - 1)X_2\,\be_2, \qquad
  \sigma_{11} = 2(\Pone + \Ptwo)\Bigl(\lambda^2 - \frac1{\lambda^2}\Bigr), \qquad
  \sigma_{33} = 2\Bigl(1 - \frac1{\lambda^2}\Bigr)\bigl(\Pone + \lambda^2\Ptwo\bigr),
\end{equation}
\begin{equation}
  p = \frac{\sigma_{11} + \sigma_{33}}{3}, \qquad
  P_{11} = \frac{\sigma_{11}}{\lambda}, \qquad
  \sigma_{vm} = \sqrt{\sigma_{11}^2 + \sigma_{33}^2 - \sigma_{11}\sigma_{33}} .
\end{equation}
\begin{align}
  \text{NH:}\quad & \sigma_{11} = \mu\Bigl(\lambda^2 - \frac1{\lambda^2}\Bigr), &
  \sigma_{33} &= \mu\Bigl(1 - \frac1{\lambda^2}\Bigr), &
  p &= \frac{\mu}{3}\Bigl(\lambda^2 + 1 - \frac{2}{\lambda^2}\Bigr), \\
  \text{MR:}\quad & \sigma_{11} = \mu\Bigl(\lambda^2 - \frac1{\lambda^2}\Bigr), &
  \sigma_{33} &= 2\bigl(c_1 + c_2\lambda^2\bigr)\Bigl(1 - \frac1{\lambda^2}\Bigr), &
  p &= \frac{\sigma_{11} + \sigma_{33}}{3}.
\end{align}
To realise this in the code (Dirichlet conditions are applied to the full
displacement vector), prescribe the affine $\bu$ above on the two end
faces $X_1 = 0$ and $X_1 = L$ and leave $X_2 = 0, H$ free; alternatively
prescribe $\bu$ on $X_1 = 0$ and the nominal traction
$\bt_R = P_{11}\,\be_1$ on $X_1 = L$.  With Taylor--Hood spaces the affine
displacement and the constant pressure lie in the discrete spaces, so the
mixed formulation reproduces $\bu$, the constant $p$ above and
$\sigma_{vm}$ to solver tolerance.  The example meshes are generated by Gmsh from \code{apps/mesh/*.geo} with
named physical groups (\code{left}, \code{right}, \code{top}, \code{bottom}),
and the input files name the faces rather than numbering them.

\subsection{Plane-strain simple shear}
$\bu = \gamma X_2\,\be_1$.  With the displacement prescribed on the whole
boundary (the setting of the patch test in \code{tests/test\_mixed.cpp})
the incompressible pressure is determined only up to a constant, because
adding a constant to $p$ adds $c\int J\bF^{-T}:\Grad\bw\ud V = c\int_{\partial\Omega_R}\bw\cdot J\bF^{-T}\bN\ud A = 0$
to the residual for every admissible $\bw$ (Piola identity).  The
checkable quantities are then the deviatoric stress $\bsig'$, the shear
stress $\sigma_{12} = 2(\Pone + \Ptwo)\gamma$ and the normal-stress
differences of Section \ref{sec:homogeneous}; at finite $\kappa$ the
affine field is exactly isochoric, so $p = \kappa(J - 1) = 0$ and
$\bsig = \bsig'$.  If instead only the top and bottom faces are constrained
and the lateral faces are left free, simple shear is \emph{not} the exact
solution (the lateral faces would have to carry $P_{11} = \sigma_{11} - \gamma\sigma_{12}$
and $P_{21} = \sigma_{12} - \gamma\sigma_{22}$); the neo-Hookean model is the
exception, with $P_{11} = 0$ but $P_{21} = \mu\gamma \ne 0$, so even there
free lateral faces perturb the state near the edges.

\subsection{Plane stress}
\label{sec:planestress}

For a thin sheet loaded in its plane the appropriate two-dimensional
idealisation is plane stress, $\sigma_{i3} = 0$.  The thickness stretch
$\lambda_3 = F_{33}$ is then not prescribed but set by the material.  For an
incompressible material $\lambda_3 = 1/\det\bF_{2D}$, and $\sigma_{33} = 0$
fixes the mean stress as $p = -\sigma'_{33}$, so
\begin{equation}
  \sigma_{\alpha\beta} = \sigma'_{\alpha\beta} - \sigma'_{33}\,\delta_{\alpha\beta},
  \qquad
  \bP = \bP_{\mathrm{iso}} + p\,\bF^{-T}, \quad p = -\lambda_3 P_{\mathrm{iso},33} .
\end{equation}
The constraint and its multiplier are eliminated pointwise: the sheet
problem is a displacement problem in the two in-plane components, with no
pressure unknown, no inf-sup condition and no locking.  For a compressible
material $\lambda_3$ solves $P_{33}(\bF_{2D}, \lambda_3) = 0$ at every point.
The states of Section \ref{sec:homogeneous} are all plane-stress states, so
they become two-dimensional problems: the uniaxial sheet with
$\bu = (\lambda - 1)X_1\,\be_1 + (\lambda^{-1/2} - 1)X_2\,\be_2$ on the end
faces and the lateral faces free, and equibiaxial tension, pure shear and
thin-sheet simple shear ($\sigma_{33} = 0$, Section 4.4) with the affine
displacement on the whole boundary, which is well posed here because $p$
is determined pointwise (contrast Section 5.2).  The code selects it with
\code{plane: stress} (2D, displacement formulation); the material is
wrapped in an adapter that returns $\bP$ at the thickness stretch, so the
tangent obtained by dual numbers is the consistent plane-stress tangent
(for a compressible material the converged root is differentiated
implicitly).  The thickness stretch is the output field
\code{thickness\_stretch}, and \code{jacobian} is the true $J$ (one for an
incompressible sheet).  Since each load step starts with the boundary
displaced and the interior lagging, shear-dominated sheet problems need
increments that shear a boundary layer of elements by no more than a few
percent, as the energy grows like $(\det\bF_{2D})^{-2}$.

\subsection{Parameter correspondence}

\begin{table}[htbp]
\centering
\small
\begin{tabular}{@{}l l >{\raggedright\arraybackslash}p{4.4cm}@{}}
\toprule
Model & \code{material} keys & Notes \\
\midrule
Neo-Hookean & \code{model: iso\_neo\_hookean}, \code{mu} & or \code{E, nu} with $\mu = E/(2(1+\nu))$ \\
Mooney--Rivlin & \code{model: mooney\_rivlin}, \code{c1}, \code{c2} & $\mu = 2(c_1 + c_2)$ \\
Yeoh & \code{model: yeoh}, \code{c10}, \code{c20}, \code{c30} & $\mu = 2C_{10}$; \code{c20}, \code{c30} default to 0 \\
Gent & \code{model: gent}, \code{mu}, \code{Jm} & \\
Arruda--Boyce & \code{model: arruda\_boyce}, \code{mu}, \code{N} & \code{mu} is $nkT$, $\mu_0$ from \eqref{eq:ab_mu0} \\
Ogden & \code{model: ogden}, \code{mu\_r}, \code{alpha\_r} & lists of up to six terms \\
Incompressible & \code{incompressible: true} or \code{nu: 0.5} & needs \code{formulation: mixed} \\
Near-incompressible & \code{kappa:} $\kappa$ or \code{nu:} $\nu < 0.5$ & $p = \kappa(J-1)$, either formulation \\
\bottomrule
\end{tabular}
\caption{Mapping between the models above and the material block of the input file.}
\end{table}

\noindent
The \code{pressure} output field is $p = \tfrac13\tr\bsig$, tensile
positive; the first Piola--Kirchhoff stress assembled by the mixed kernel is
$\bP = \bP_{\mathrm{iso}} + p\,J\bF^{-T}$.  A compressive hydrostatic state
therefore has $p < 0$.

\subsection{Verification inputs}
\label{sec:verification}

{\sloppy
The inputs \code{apps/input/finite\_elasticity/homogeneous/plane\_strain\_<model>.yaml} (unit
square \code{apps/mesh/square.msh}, plane-strain extension to $\lambda = 2$
as in Section 5.1),
\code{plane\_stress\_<model>.yaml} (the same sheet in plane stress: uniaxial
tension to $\lambda = 2$, thickness stretch $\lambda^{-1/2}$, no pressure
unknown) and \code{uniaxial\_<model>.yaml} (unit cube \code{apps/mesh/cube.msh},
uniaxial tension to $\lambda = 2$, the affine displacement of Section 4.1
prescribed on the faces \code{left} and \code{right}, the four lateral faces
free) exist for all six models, together with equibiaxial and pure-shear
cubes (\code{equibiaxial\_<model>.yaml}, \code{pure\_shear\_<model>.yaml};
Sections 4.2 and 4.3, the faces $Z = 0, 1$ free) and the uniaxial symmetry
model with rollers.  Their Dirichlet entries write the affine data
$\bu = \bm{G}\bX$, $\bm{G} = \operatorname{diag}(\lambda_a - 1)$, as
expression strings (e.g.\ \code{["x", "-0.5*y"]}).
\code{apps/homogeneous\_compare.py} runs them and prints the probed
displacement, pressure, von Mises stress, $J$, and the components of
$\bsig$ and $\bP$ next to the closed forms of Sections 4.1 and 5.1
(\code{make homogeneous}); all agree to about $10^{-13}$ with quadratic
Newton convergence, in the nodal and the element presentations alike.  \code{tests/test\_homogeneous}
(part of \code{make check}) repeats the comparison programmatically: at the
material point, the code's Cauchy stress with the analytic mean stress
against the principal stresses of Sections 4.1--4.3 and 5.1 and the
deviatoric stress of simple shear (Section 4.4); as finite element
solutions of the mixed formulation in 2D and 3D, with every quadrature
quantity checked in its nodal (both projections), element-average and
quadrature-point presentations; and as plane-stress
solutions of the sheet states uniaxial, equibiaxial, pure shear and simple
shear (Section \ref{sec:planestress}), including the Mooney--Rivlin
$\sigma_{22} = -2c_2\gamma^2$ of simple shear.\par}

%============================================================================
\small\begin{thebibliography}{9}
%============================================================================
\bibitem{ogden} R.\,W. Ogden, \emph{Non-Linear Elastic Deformations}, Ellis Horwood, 1984 (Dover, 1997).
\bibitem{holzapfel} G.\,A. Holzapfel, \emph{Nonlinear Solid Mechanics: A Continuum Approach for Engineering}, Wiley, 2000.
\bibitem{treloar} L.\,R.\,G. Treloar, \emph{The Physics of Rubber Elasticity}, 3rd ed., Oxford University Press, 1975.
\bibitem{mooney} M. Mooney, A theory of large elastic deformation, \emph{J. Appl. Phys.} 11 (1940) 582--592.
\bibitem{rivlin} R.\,S. Rivlin, Large elastic deformations of isotropic materials IV, \emph{Phil. Trans. R. Soc. A} 241 (1948) 379--397.
\bibitem{yeoh} O.\,H. Yeoh, Some forms of the strain energy function for rubber, \emph{Rubber Chem. Technol.} 66 (1993) 754--771.
\bibitem{gent} A.\,N. Gent, A new constitutive relation for rubber, \emph{Rubber Chem. Technol.} 69 (1996) 59--61.
\bibitem{ab} E.\,M. Arruda, M.\,C. Boyce, A three-dimensional constitutive model for the large stretch behavior of rubber elastic materials, \emph{J. Mech. Phys. Solids} 41 (1993) 389--412.
\bibitem{ogden72} R.\,W. Ogden, Large deformation isotropic elasticity: on the correlation of theory and experiment for incompressible rubberlike solids, \emph{Proc. R. Soc. A} 326 (1972) 565--584.
\end{thebibliography}

\end{document}
